\documentclass[10pt,twocolumn]{article}
\usepackage[T1]{fontenc}
\usepackage[utf8]{inputenc}
\usepackage{lmodern}
\usepackage[margin=1.8cm,top=2.4cm,bottom=2cm,columnsep=0.6cm]{geometry}
\usepackage{fancyhdr}
\usepackage{titlesec}
\usepackage{booktabs}
\usepackage{amsmath,amssymb,amsfonts}
\usepackage{microtype}
\usepackage{xcolor}
\usepackage{mdframed}
\usepackage{parskip}
\usepackage{enumitem}
\usepackage{hyperref}

\definecolor{rulered}{RGB}{180,30,30}
\definecolor{boxgray}{RGB}{245,245,245}
\definecolor{keywordgray}{RGB}{100,100,100}

\pagestyle{fancy}
\fancyhf{}
\fancyhead[L]{\textsc{\small Claude Mag}}
\fancyhead[C]{\small\textit{Machine Learning Research Digest}}
\fancyhead[R]{\small 2026-11-04}
\fancyfoot[C]{\small\thepage}
\renewcommand{\headrulewidth}{0.8pt}

\titleformat{\section}{\normalsize\bfseries\color{rulered}}{}{0pt}{}%
  [\vspace{-4pt}\color{rulered}\rule{\columnwidth}{0.4pt}]
\titlespacing{\section}{0pt}{8pt}{4pt}

\hypersetup{colorlinks=true,urlcolor=rulered,linkcolor=black}

\begin{document}
\setlength{\parindent}{0pt}
\setlength{\parskip}{4pt}

{\small\color{keywordgray}\textbf{Keywords:}
  continual learning \textbullet{}
  retrieval-augmented agents \textbullet{}
  tutorial video understanding \textbullet{}
  nonparametric adaptation}\par\vspace{4pt}

{\LARGE\bfseries\raggedright
  Can LLM Agents Learn by Watching Tutorial Videos?}\par\vspace{4pt}

{\large\itshape\raggedright
  Tutorial Agentic Learning Turns Video Procedures into Frozen-LLM Memory
  with No Weight Updates}\par\vspace{6pt}

{\small\textbf{By} Haoyu Zhao, Zafeirios Fountas, Jun Wang}\par\vspace{2pt}

{\small\color{keywordgray}arXiv:2611.01460 \textbullet{}
  cs.LG / cs.AI / cs.CL / cs.CV \textbullet{} Submitted 1 Nov 2026}
\par\vspace{2pt}\noindent{\color{rulered}\rule{\columnwidth}{0.6pt}}\vspace{6pt}

Tutorial Agentic Learning (TAL) enables a frozen large language model to
acquire procedural knowledge from instructional videos at inference time,
without any weight updates. Tested on VidLearn---a new multimodal
benchmark of 72 chess tutorial videos spanning 12 openings---TAL
performance scales with the number of videos consumed, demonstrating
genuine procedural learning from visual and language content.

\begin{mdframed}[backgroundcolor=boxgray,linecolor=rulered,linewidth=1pt,
    innerleftmargin=6pt,innerrightmargin=6pt,
    innertopmargin=6pt,innerbottommargin=6pt]
\textbf{\small AT A GLANCE}\par\vspace{4pt}
\begin{description}[leftmargin=0pt,labelsep=4pt,itemsep=2pt,parsep=0pt,
    font=\small\bfseries]
  \item[Problem:] LLM agents require expensive retraining to acquire new
    procedural skills, blocking continual learning at deployment.
  \item[Why it matters:] Tutorial videos are abundant and contain rich
    procedural knowledge; efficiently converting them to agent capability
    would eliminate per-skill fine-tuning cycles.
  \item[Method:] Videos are segmented into timestamped procedural entries
    and stored in a retrieval-augmented memory bank; at inference time,
    top-$k$ entries provide in-context guidance to a frozen LLM.
  \item[Result:] TAL outperforms the same LLM without video context;
    accuracy scales with number of videos ingested on the VidLearn
    benchmark.
  \item[Evidence:] VidLearn: 72 chess tutorial videos, 12 openings,
    multimodal QA; dense retrieval ablations over retrieval count $k$.
\end{description}
\end{mdframed}

\section{The Big Picture}

The boundary between training and inference has been treated as fixed in
standard machine learning: an agent learns during training and exploits
what it learned during deployment. But in real-world settings, new
procedural knowledge becomes available continuously---software tutorials,
instructional videos, how-to guides---and re-initiating a training run
for each new skill is impractical.

Tutorial Agentic Learning (TAL) reframes this problem: rather than
modifying weights, it treats instructional videos as a persistent,
queryable memory that augments a frozen agent at inference time. The
approach is inspired by retrieval-augmented generation (RAG) for
factual QA, but extends the paradigm to procedural, visually grounded
knowledge embedded in sequential video content.

\section{Why Existing Approaches Fall Short}

\textbf{Fine-tuning on instructional data} requires labeled examples and
gradient computation, is prone to catastrophic forgetting of previously
learned skills, and must be repeated for each new skill domain.

\textbf{In-context learning from text summaries} loses the visual content
of videos and fails to capture the sequential, state-action structure of
procedural knowledge, particularly in domains like chess or software
manipulation where board/screen state is essential.

\textbf{Vision-language models applied to full video} at inference time
are computationally prohibitive: processing 72 tutorial videos at every
query step would require multiple full video passes per query.

\section{Core Mechanism}

TAL operates in two phases. During \textbf{offline indexing}, each video
is segmented into timestamped steps using a vision-language model (VLM).
Each step is converted into a structured procedural entry:
\[
  e_i = (\text{state}_i,\; \text{action}_i,\; \text{rationale}_i)
\]
where $\text{state}_i$ is a visual embedding of the board/screen,
$\text{action}_i$ is the described move or operation, and
$\text{rationale}_i$ is the narrated explanation. All entries are stored
in a FAISS-indexed memory bank.

During \textbf{inference-time retrieval}, given a query $q$, the top-$k$
entries are retrieved via dense similarity search:
\[
  \mathcal{E}^* = \arg\!\top_k \bigl\{\cos(\phi(q),\, \phi(e_i))\bigr\}_{i=1}^{N}
\]
where $\phi(\cdot)$ is a frozen CLIP encoder. The retrieved entries are
formatted as structured in-context examples and provided to the frozen
LLM for final answer generation.

\section{The VidLearn Benchmark}

VidLearn is a new multimodal QA benchmark constructed from 72 publicly
available chess tutorial videos spanning 12 chess openings (Sicilian
Defense, King's Indian, French Defense, and 9 others). Each video is
annotated with key board positions, move-by-move explanations, and
multiple-choice and free-form questions. The benchmark tests whether an
agent can correctly identify optimal continuations from visually presented
board states, drawing on knowledge explicitly taught in the videos.

The 12 openings are partitioned into train (8) and test (4) splits to
evaluate cross-opening transfer.

\section{Learning or Inference Procedure}

\textbf{Offline (once per video set):}
\begin{enumerate}[itemsep=2pt,parsep=0pt]
  \item Segment each video into $T_v$ timestamped steps using a VLM.
  \item Embed each step as a (state, action, rationale) triple.
  \item Encode each triple with CLIP and add to the FAISS index.
\end{enumerate}

\textbf{Inference (per query):}
\begin{enumerate}[itemsep=2pt,parsep=0pt]
  \item Encode the query board image and question text with CLIP.
  \item Retrieve top-$k$ entries from the memory bank.
  \item Format as in-context examples and pass to the frozen LLM.
  \item Return the LLM's generated answer.
\end{enumerate}

No gradient computation occurs at any stage of inference.

\section{What the Guarantee Says}

TAL provides no formal learning-theoretic guarantee. The empirical claim
is a scaling law: accuracy on VidLearn test questions increases as the
number of ingested training-opening videos grows, leveling off around
50 videos per opening. This suggests that the memory bank successfully
captures the unique procedural patterns in a given opening up to a
saturation point. No guarantee is provided for out-of-distribution
openings, though cross-opening transfer is observed for structurally
related sub-variants.

\section{Experimental Findings}

\begin{itemize}[itemsep=2pt,parsep=0pt]
  \item \textbf{Accuracy vs.\ baseline:} TAL with 72 videos achieves
    substantially higher accuracy than the same LLM without video context
    on VidLearn multiple-choice questions.
  \item \textbf{Scaling:} Accuracy increases monotonically with the
    number of ingested videos, with diminishing returns above 50 videos
    per opening.
  \item \textbf{Cross-opening transfer:} Positive transfer is observed
    for Sicilian Defense sub-variants, with 4.1\% accuracy improvement on
    unseen variations when related-opening videos are present.
  \item \textbf{Retrieval ablation:} Performance improves with $k$ up to
    $k = 5$, then plateaus, suggesting diminishing returns from additional
    in-context examples beyond 5.
\end{itemize}

\section{Ablations and Interpretation}

\textbf{Without visual state embeddings:} Replacing CLIP visual embeddings
with text-only embeddings of move descriptions reduces accuracy by 7.3\%,
confirming that the board-state visual signal is essential for precise
retrieval.

\textbf{Without the rationale field:} Removing the rationale from
in-context entries reduces accuracy by 4.9\%, indicating that the LLM
uses explanatory text from videos to reason about move motivations.

\textbf{Approximate vs.\ exact search:} Using exact L2 search instead of
approximate FAISS search improves accuracy by less than 0.5\%, confirming
that approximate retrieval is sufficient and scalable.

\textbf{Impact of VLM quality for segmentation:} Replacing the full VLM
with a smaller captioning model for step segmentation reduces accuracy by
6.1\%, highlighting that the quality of initial video parsing matters.

\section*{Reference}

Haoyu Zhao, Zafeirios Fountas, Jun Wang.
\textbf{Can Large Language Model Agents Learn by Watching Tutorial
Videos?}
arXiv:2611.01460. 2026.
\url{https://arxiv.org/abs/2611.01460}

\end{document}
